\newpage

\section{Consuntivo di periodo}
\subsection{Introduzione}
La presente sezione raccoglie, per ciascuno sprint, il confronto tra quanto pianificato e quanto effettivamente svolto dal gruppo \textbf{Git Happens}, in termini di attività, ore produttive e costi. \\
\\
Al termine di ogni sprint bisettimanale il Responsabile di turno compila il consuntivo durante la riunione interna di chiusura, confrontando le ore preventivate con quelle rendicontate da ciascun componente e aggiornando la tabella \textit{Totale ore rendicontate} della Sezione~4.4. Ogni sprint riporta inoltre i rischi attesi e quelli effettivamente incontrati, con riferimento ai codici censiti nella Sezione~2, e si chiude con una retrospettiva da cui derivano le migliorie da attuare nello sprint successivo, secondo il ciclo \textbf{PDCA} descritto nella Sezione~3. \\
\\
Uno scostamento superiore al 15\% tra ore preventivate e consuntivate costituisce indicatore di rilevamento del rischio \textbf{RO-01} e comporta l'attivazione delle relative misure di mitigazione.

\subsection{Verso la Requirements and Technology Baseline}
Gli sprint di questo periodo sono orientati alla revisione \textbf{RTB}, prevista entro il 15/01/2027: consolidamento dei processi e della documentazione di progetto, analisi dei requisiti e realizzazione del \textit{Proof of Concept} (PoC) a validazione dello stack tecnologico.

\subsubsection{Sprint 1}

\paragraph{Informazioni generali}~\\[0.15cm]
\begin{itemize}
    \item \textbf{Periodo:} dal 10/11/2026 al 24/11/2026;
    \item \textbf{Obiettivi:} predisposizione dell'infrastruttura di lavoro (account cloud condivisi, repository del PoC, pipeline CI, spazi Jira), avvio della documentazione di progetto (\textit{Norme di Progetto}, \textit{Piano di Progetto}, \textit{Analisi dei Requisiti}), prima esplorazione dello stack tecnologico e chiarimento dei dubbi sul capitolato con il proponente \textbf{Halue S.r.l. SB} (incontri dell'11/11/2026 e del 18/11/2026).
\end{itemize}

\noindent
I ruoli assegnati per lo sprint (\texttt{VI\_2026-11-10\_D01}) sono i seguenti:

\vspace{2mm}
\renewcommand{\arraystretch}{1.3}
\noindent
\begin{tabularx}{\textwidth}{|l X|}
\hline
\textbf{Ruolo} & \textbf{Componente} \\
\hline
Responsabile & Cristiano Antoniazzo \\
Amministratore & Antonio Hai Xu Tang \\
Progettista & Filippo Lissandrin \\
Analista & Gabriele Gambarro, Diana Grosselle, Gianmaria Grinovero \\
Verificatore & Tommaso Caon \\
\hline
\end{tabularx}

\paragraph{Attività svolte}~\\[0.15cm]
{\small
\renewcommand{\arraystretch}{1.3}
\noindent
\begin{tabularx}{\textwidth}{|>{\raggedright\arraybackslash}X >{\raggedright\arraybackslash}p{2.6cm} >{\raggedright\arraybackslash}p{1.6cm} c|}
\hline
\textbf{Attività} & \textbf{Assegnatario} & \textbf{Ticket} & \textbf{Stato} \\
\hline
Registrazione degli account condivisi (Vercel, Supabase, Shopify, Google OAuth) e foglio dei limiti dei piani gratuiti & C. Antoniazzo & -- & Completata \\
Repository del PoC, pipeline CI di base e protezione del branch \texttt{main} & A. Tang & -- & Completata \\
Configurazione della pipeline Vercel Pro e integrazione con GitHub & F. Lissandrin, A. Tang & C8-11 & Completata \\
Space Jira per C8 e documentazione delle tecnologie & F. Lissandrin & -- & Completata \\
Collegamento dello Space Jira a Shopify & F. Lissandrin & -- & Non svolta \\
Esperimento didattico sullo stack e video dimostrativo per il team & F. Lissandrin & -- & Completata \\
Struttura del Glossario e script di automazione & G. Gambarro, A. Tang, F. Lissandrin, T. Caon & SWE-91, SWE-106 & Completata \\
Stesura e aggiornamento delle \textit{Norme di Progetto} & G. Gambarro, G. Grinovero & SWE-84, SWE-104 & In corso \\
Stesura del \textit{Piano di Progetto} & D. Grosselle, G. Grinovero, C. Antoniazzo & SWE-86, SWE-102, SWE-105 & In corso \\
Avvio dell'\textit{Analisi dei Requisiti} & D. Grosselle & SWE-103 & In corso \\
Verbali interni (10/11, 17/11) ed esterni (11/11, 18/11) & G. Grinovero, D. Grosselle, G. Gambarro & SWE-92, SWE-93, SWE-101, SWE-107 & Completata \\
\hline
\end{tabularx}}

\paragraph{Rischi attesi}~\\[0.15cm]
\begin{itemize}
    \item \textbf{RT-01} -- Curva di apprendimento dello stack: primo contatto del gruppo con Next.js, Drizzle ORM, Supabase e Vercel;
    \item \textbf{RT-02} -- Limiti dei piani gratuiti: configurazione dei primi account cloud e censimento delle quote;
    \item \textbf{RO-01} -- Sottostima delle ore: primo sprint con stime formali, senza dati storici di riferimento;
    \item \textbf{RO-02} -- Concentrazione delle competenze: l'esplorazione dello stack è affidata a pochi componenti;
    \item \textbf{RC-01} -- Ambiguità dei requisiti: diversi dubbi interpretativi sul capitolato ancora aperti (token Shopify, organizzazioni, cron, piano Vercel).
\end{itemize}

\paragraph{Preventivo}~\\[0.15cm]
La stima delle ore è stata effettuata all'avvio dello sprint sulla base delle attività assegnate a ciascun componente e valorizzata con le tariffe orarie della Sezione~4.2.

\vspace{2mm}
\renewcommand{\arraystretch}{1.3}
\noindent
\begin{tabularx}{\textwidth}{|X l c c|}
\hline
\textbf{Componente} & \textbf{Ruolo} & \textbf{Ore} & \textbf{Costo (\texteuro{})} \\
\hline
Cristiano Antoniazzo & Responsabile & 5 & 150,00 \\
Antonio Hai Xu Tang & Amministratore & 4 & 80,00 \\
Filippo Lissandrin & Progettista & 5 & 125,00 \\
Gabriele Gambarro & Analista & 5 & 125,00 \\
Diana Grosselle & Analista & 6 & 150,00 \\
Gianmaria Grinovero & Analista & 4 & 100,00 \\
Tommaso Caon & Verificatore & 5 & 75,00 \\
\hline
\textbf{TOTALE} & & \textbf{34} & \textbf{805,00} \\
\hline
\end{tabularx}

\paragraph{Consuntivo}~\\[0.15cm]
Le ore effettive sono ricavate dal registro ore giornaliero compilato dai componenti durante lo sprint.

\vspace{2mm}
\renewcommand{\arraystretch}{1.3}
\noindent
\begin{tabularx}{\textwidth}{|X l c c|}
\hline
\textbf{Componente} & \textbf{Ruolo} & \textbf{Ore} & \textbf{Costo (\texteuro{})} \\
\hline
Cristiano Antoniazzo & Responsabile & 4,5 & 135,00 \\
Antonio Hai Xu Tang & Amministratore & 2 & 40,00 \\
Filippo Lissandrin & Progettista & 4 & 100,00 \\
Gabriele Gambarro & Analista & 4 & 100,00 \\
Diana Grosselle & Analista & 7 & 175,00 \\
Gianmaria Grinovero & Analista & 4,5 & 112,50 \\
Tommaso Caon & Verificatore & 4,5 & 67,50 \\
\hline
\textbf{TOTALE} & & \textbf{30,5} & \textbf{730,00} \\
\hline
\end{tabularx}

\vspace{3mm}
\noindent
Il confronto per ruolo tra preventivo e consuntivo è riportato nella tabella seguente.

\vspace{2mm}
\renewcommand{\arraystretch}{1.3}
\noindent
\begin{tabularx}{\textwidth}{|X c c c c c|}
\hline
\textbf{Ruolo} & \textbf{Ore prev.} & \textbf{Ore eff.} & \textbf{Scost.} & \textbf{Costo prev. (\texteuro{})} & \textbf{Costo eff. (\texteuro{})} \\
\hline
Responsabile & 5 & 4,5 & $-$0,5 & 150,00 & 135,00 \\
Amministratore & 4 & 2 & $-$2 & 80,00 & 40,00 \\
Progettista & 5 & 4 & $-$1 & 125,00 & 100,00 \\
Analista & 15 & 15,5 & +0,5 & 375,00 & 387,50 \\
Verificatore & 5 & 4,5 & $-$0,5 & 75,00 & 67,50 \\
\hline
\textbf{TOTALE} & \textbf{34} & \textbf{30,5} & \textbf{$-$3,5} & \textbf{805,00} & \textbf{730,00} \\
\hline
\end{tabularx}

\vspace{3mm}
\noindent
Lo sprint si chiude con 3,5 ore in meno rispetto al preventivo ($-$10,3\%) e un risparmio di \texteuro{} 75,00. Lo scostamento resta entro la soglia di allarme del 15\% (RO-01). Lo scarto maggiore riguarda l'Amministratore ($-$2 ore): le attività infrastrutturali (protezione del branch \texttt{main}, sistemazione di Vercel e del Glossario) hanno richiesto meno tempo del previsto. Il minor impegno del Progettista dipende dal collegamento a Shopify, che non è stato svolto ed è riportato allo sprint successivo. L'unico ruolo oltre la stima è l'Analista (+0,5 ore), che ha lavorato in parallelo su \textit{Norme di Progetto}, \textit{Piano di Progetto}, \textit{Analisi dei Requisiti} e verbali. \\
\\
Al termine dello sprint risultano rendicontate 30,5 ore su 651, per un costo di \texteuro{} 730,00 su \texteuro{} 12.845,00. Restano disponibili 620,5 ore e un budget residuo di \texteuro{} 12.115,00.

\paragraph{Retrospettiva}~\\[0.15cm]
% TODO: aspetti positivi e negativi emersi nella riunione di chiusura
\textit{Da compilare nella riunione di chiusura dello sprint.}

\paragraph{Rischi incontrati}~\\[0.15cm]
% TODO
\textit{Da compilare nella riunione di chiusura dello sprint.}

\paragraph{Strategia di mitigazione adottata}~\\[0.15cm]
% TODO
\textit{Da compilare nella riunione di chiusura dello sprint.}

\paragraph{Migliorie da attuare}~\\[0.15cm]
% TODO
\textit{Da compilare nella riunione di chiusura dello sprint.}

